\color{unidarkgreen}{\subsection{Дистанционная защита (ДЗ)}}
\color{black}

\begin{enumerate}[label=\arabic{section}.\arabic{subsection}.\arabic*,  align=left, labelsep=4pt, leftmargin=0pt, itemindent=57pt]

\item
ИЭУ содержит четыре ступени ДЗ от междуфазных замыканий и четыре ступени ДЗ от замыканий
на землю. Дистанционная  защита  предназначена  для  селективного  отключения  поврежденного  элемента
электрической  сети  при  всех видах  коротких  замыканий (КЗ),  реагируя  на снижение измеряемого
сопротивления в защищаемой зоне.

\item 
ДЗ селективно срабатывает:
\begin{itemize}
    \item при всех видах КЗ; 
    \item при переходе внешнего КЗ во внутреннее; 
    \item при КЗ, сопровождающемся снижением измеряемого напряжения до нуля.
\end{itemize}

\item 
Обеспечивается несрабатывание защиты:
\begin{itemize} 
\item при качаниях и асинхронном ходе; 
\item при неисправностях цепей напряжения;
\item при внешних КЗ с насыщением трансформаторов тока (при соблюдении заявленных производителем требований к ТТ; 
\item при бросках тока намагничивания без КЗ;
\item при внешних КЗ, даже если измеряемое напряжение снижается до нуля.
\end{itemize}

\item Каждая из ступеней ДЗ от междуфазных КЗ включает в себя три реле сопротивления (РС), использующие
для замера сопротивления  комплексные значения линейных напряжений и соответствующих им линейных
токов:
% ФОРМУЛЫ
\begin{align}
\vec Z_{AB} =  \frac{\vec U_A - \vec U_B}{\vec I_A - \vec I_B}\label{eq:groupZ1}, \\[10pt]
\vec Z_{BC} =  \frac{\vec U_B - \vec U_C}{\vec I_B - \vec I_C}\label{eq:groupZ2}, \\[10pt]
\vec Z_{CA} =  \frac{\vec U_C - \vec U_A}{\vec I_C - \vec I_A}\label{eq:groupZ3}, 
\end{align}
\begin{eqexpl}[20mm]
    \item{$ \vec U_{A}, \vec U_{B}, \vec U_{C} $} комплексные значения напряжений фаз А, В, С от ТН 1СШ;
	\vspace{2pt}
    \item{$ \vec I_{A}, \vec I_{B}, \vec I_{C} $} комплексные значения токов фаз А, В, С.
\end{eqexpl}
\vspace{4pt}

\item 
Комплексные значения сопротивлений по контуру «фаза-земля» рассчитываются с учетом компенсации
тока нулевой последовательности:
% ФОРМУЛЫ
\begin{align}
\vec Z_{A} =  \frac{\vec U_A }{\vec I_A + \vec K \cdot  3 \vec I_0}\label{eq:groupZ11}, \\[10pt]
\vec Z_{B} =  \frac{\vec U_B }{\vec I_B + \vec K \cdot  3 \vec I_0}\label{eq:groupZ12}, \\[10pt]
\vec Z_{C} =  \frac{\vec U_C }{\vec I_C + \vec K \cdot  3 \vec I_0}\label{eq:groupZ13}, 
\end{align}
\begin{eqexpl}[5mm]
    \item{$ 3 \vec I_{0} $}  комплексное  значение  утроенного  тока  нулевой  последовательности  защищаемого объекта (расчетное значение);
    \item{$ \vec K $}  комплексный коэффициент компенсации тока нулевой последовательности защищаемого объекта.
\end{eqexpl}
\vspace{4pt}

\item 
В ИЭУ коэффициент  $\vec K$  задается двумя уставками, \textbf{«KR - коэффициент компенсации тока
3I0 по активному сопротивлению»}  и  \textbf{«KХ  -  коэффициент  компенсации тока  3I0 по реактивному сопротивлению»}.  

\item 
Коэффициенты компенсации рассчитываются по следующей формуле:
% ФОРМУЛЫ
\begin{align}
\vec K = \frac{1}{3} \cdot \frac{\vec Z_{0 \textsl{ уд}} - \vec Z_{1 \textsl{ уд}}}{\vec Z_{1 \textsl{ уд}}}\label{eq:groupK1},
\end{align}

\begin{eqexpl}
\item{$ \vec Z_{0 \textsl{ уд}} $}  комплексное удельное сопротивление нулевой последовательности;
\item{$ \vec Z_{1 \textsl{ уд}} $}  комплексное удельное сопротивление прямой последовательности.
\end{eqexpl}
\vspace{4pt}

\item 
Характеристики срабатывания ненаправленного РС для междуфазных контуров задаются независимо
для каждой ступени уставками:
\begin{itemize} 
\item \textbf{«Rср ДЗ-ФФ-1(2..4)»} – активное сопротивление срабатывания РС; 
\item \textbf{«Хср ДЗ-ФФ-1(2..4)»} – реактивное сопротивление срабатывания РС;
\item \textbf{«Ф1 ДЗ-ФФ-1(2..4)»} – угол линии электропередач; 
\end{itemize}

\item 
Характеристики срабатывания ненаправленного РС для фазных контуров задаются независимо
для каждой ступени уставками:
\begin{itemize} 
\item \textbf{«Rср ДЗ-ФЗ-1(2..4)»} – активное сопротивление срабатывания РС; 
\item \textbf{«Хср ДЗ-ФЗ-1(2..4)»} – реактивное сопротивление срабатывания РС;
\item \textbf{«Ф1 ДЗ-ФЗ-1(2..4)»} – угол линии электропередач; 
\end{itemize}

\item 
Для предотвращения неселективного срабатывания первой и третьей ступеней при внешних
повреждениях через  переходное  сопротивление  предусмотрена  возможность  отсечения области
характеристики срабатывания,  задаваемой  углом  $φ_4$.  Область  отсекается  отрезком, исходящей  под  углом  $φ_4$,
из точки на верхней части характеристики, соответствующей углу линии электропередачи $φ_1$. Угол $φ_4$
задается уставками \textbf{«Ф4 ДЗ-ФФ-1(3)»} и \textbf{«Ф4 ДЗ-ФЗ-1(3)»}. 

\item 
Вид характеристики срабатывания ненаправленного РС представлен на рисунке \ref{pict_dz:char1}. 

\begin{figure}[H]
\centering
\includegraphics[width=0.4\textwidth,height=0.4\textheight,keepaspectratio]{\fbpath/661.0031 ДЗ/img/_char1.pdf}
\captionof{figure}{ Характеристика срабатывания ненаправленного РС}\label{pict_dz:char1}
\end{figure}

\item 
Параметры реле сопротивления приведены в таблице \ref{desc_dz:tbl1}

\small
\setlength{\extrarowheight}{0.05cm} % высота ячеек в таблице
\begin{longtable}
{|>{\centering\arraybackslash}m{15.2cm}|>{\centering\arraybackslash}m{1.7cm}|}
\caption{Параметры РС\hfill\vspace{-0.5\baselineskip}}\label{desc_dz:tbl1}\\ 
\hhline{|-|-|}
\rowcolor{gray!30}
Наименование параметра & Значение  \\ 
\hhline{|-|-|}
\endfirsthead

\caption*{\hspace{3pt}\emph{Продолжение таблицы \ref{desc_dz:tbl1}\hfill\vspace{-0.5\baselineskip}}} \\ % сделано по ГОСТ 2.105 п.6.8.7
\hhline{|-|-|}
\rowcolor{gray!30}
Наименование параметра & Значение  \\ 
\hhline{|-|-|}
\endhead
\hhline{|-|-|}

\endfoot
\endlastfoot

\raggedright Минимальное  междуфазное  напряжение,  при  котором  обеспечиваются  точностные характеристики РС, В, не менее & \centering\arraybackslash 0,5  \\  \hline
\raggedright Ток точной работы во всем диапазоне уставок при работе на угле $φ_1$, \% от номинального тока, не более  & \centering\arraybackslash 10  \\  \hline
\raggedright Собственное  время  срабатывания  РС  при  работе  на  угле  максимальной  чувствительности, трехкратном токе точной работы и скачкообразном уменьшении напряжения  на входе РС от напряжения,  соответствующего  сопротивлению  на  зажимах  РС  не  менее  $1,2\cdot(X_{\textsl{уст}} /sinφ_1 )$ до напряжения соответствующего $0,6\cdot(X_{\textsl{уст}} /sinφ_1 )$ мс, не более   & \centering\arraybackslash 25  \\  \hline
\raggedright Собственное  время  возврата  РС  при  работе  на  угле  максимальной  чувствительности, трехкратном  токе  точной  работы  и  скачкообразном  увеличении  напряжения  на  входе  РС  от напряжения соответствующего сопротивлению на зажимах РС $0,1\cdot(X_{\textsl{уст}} /sinφ_1 )$ до напряжения соответствующего $1,2\cdot(X_{\textsl{уст}} /sinφ_1 )$ (но не более 100 В для междуфазного канала и 57 В –- для фазного), мс, не более   & \centering\arraybackslash 50  \\  \hline
\raggedright Средняя основная погрешность всех РС по величине сопротивления срабатывания $R_{\textsl{уст}}$  и $X_{\textsl{уст}}$ при токе, равном $I_{\textsl{ном}}$ (или, в зависимости от уставки, меньшем токе, исходя из максимального напряжения на зажимах РС, равного 100 В), \% от уставки, не более  & \centering\arraybackslash $\pm$5  \\  \hline
\raggedright Основная абсолютная погрешность РС по углу наклона характеристики при токе, равном $I_{\textsl{ном}}$, эл. град., не более   & \centering\arraybackslash $\pm$5  \\  \hline
\raggedright Дополнительная абсолютная погрешность РС по углу наклона характеристики при изменении тока КЗ в диапазоне $(0,2 \dots 30) \cdot I_{\textsl{ном}}$, эл. град. от значения, измеренного при $I_{\textsl{ном}}$, не более   & \centering\arraybackslash $\pm$7  \\  \hline
\raggedright Дополнительная  погрешность  РС  по  сопротивлениям  $R$  и  $X$  при  изменении  температуры окружающей  среды  в  рабочем  диапазоне,  \%  от  среднего  значения,  измеренного  при нормальных климатических условиях, не более   & \centering\arraybackslash $\pm$3  \\  \hline
\raggedright Время  срабатывания  ДЗ  при  переходе  внешнего  КЗ  во  внутреннее  в  условиях  наличия насыщения ТТ мс, не более & \centering\arraybackslash 60  \\  \hline
\raggedright Коэффициент возврата, не более   & \centering\arraybackslash 1,1  \\  \hline

\end{longtable}
\setlength{\extrarowheight}{0.0cm} % возврат параметра высота ячеек в таблице
\normalsize

\item 
Выполнение функций направленного РС обеспечивается совместным использованием ненаправленного
РС и органов направленности (ОН). Орган направленности представляет собой два отрезка, исходящие из
начала координат и расположенные во втором и четвёртом квадрантах. Вид характеристики  (см. рис. \ref{pict_dz:on}) определяется
углами $φ_2$ и $φ_3$  наклона  этих  отрезков,  отсчитываемых  относительно  оси  $R$. Уставки органов направленности
являются общими для всех ступеней, задаются раздельно для контуров «фаза-фаза» (с помощью уставок \textbf{«Ф2
ОН-ФФ»} и \textbf{«Ф3 ОН-ФФ»}) и «фаза-земля» (с помощью уставок \textbf{«Ф2 ОН-ФЗ»} и \textbf{«Ф3 ОН-ФЗ»}).

\item
Углы наклона в ИЭУ задаются для прямонаправленного ОН в сторону 1СШ, а характеристика
обратнонаправленного ОН в сторону 2СШ симметрична характеристике прямонаправленного ОН с поворотом на
180 градусов.

\item
Направленность для каждой ступени ДЗ задается с помощью программных переключателей \textbf{«Направл.
ДЗ-ФФ-1(2..4)»} и \textbf{«Направл. ДЗ-ФЗ-1(2..4)»}, имеющих три положения:
\begin{enumerate}[label=\arabic*)]
\item \textit{«Ненапр»} -- ступень выполнена ненаправленной; 
\item \textit{«В 1СШ»} -- ступень выполнена прямонаправленной в сторону 1СШ; 
\item \textit{«В 2СШ»} -- ступень выполнена обратнонаправленной в сторону 2СШ. 
\end{enumerate}


\begin{figure}[H]
\centering
\includegraphics[width=0.6\textwidth,height=0.6\textheight,keepaspectratio]{\fbpath/661.0031 ДЗ/img/_on.pdf}
\captionof{figure}{Характеристики срабатывания органов направления}\label{pict_dz:on}
\end{figure}

\item
При близких трехфазных КЗ, когда все напряжения на входе РС близки к нулю, для определения
направленности в течение времени не менее 40 мс используются напряжения предаварийного режима (работа
по «памяти»). При работе РС «по памяти» при трехфазных КЗ в месте установки защиты обеспечивается
длительность сигнала срабатывания на выходе РС не менее 40 мс в диапазоне токов от 0,2$\cdot I_{\textsl{ном}}$ до 30$\cdot I_{\textsl{ном}}$.
При КЗ «за спиной» с токами до 30$\cdot I_{\textsl{ном}}$  обеспечивается отсутствие ложных срабатываний РС, в том числе при
снижении измеряемого напряжения до нуля. 

\item
Для отстройки от нагрузочного режима возможен вырез в характеристике срабатывания РС. Уставки
измерительного органа отстройки от нагрузки являются общими для всех ступеней, задаются раздельно для
контуров «фаза-фаза» и «фаза-земля».   Характеристика  срабатывания  приведена  на  рисунке \ref{pict_dz:r}.  Эта характеристика
срабатывания определяется уставками: 
\begin{itemize}
    \item \textbf{«Rнг-ФФ»} -- сопротивление выреза нагрузочного режима для контуров «фаза-фаза»;
    \item \textbf{«Фнг-ФФ»} -- угол выреза нагрузочного режима ступеней для контуров «фаза-фаза»;
	\item \textbf{«Rнг-ФЗ»} -- сопротивление выреза нагрузочного режима для контуров «фаза-земля»;
	\item \textbf{«Фнг-ФЗ»} -- угол выреза нагрузочного режима ступеней для контуров «фаза-земля».
\end{itemize}

\item
Исключаемая область в характеристике срабатывания симметрична относительно осей координат. Ввод
в работу отстройки от нагрузочного режима осуществляется с помощью программного ключа \textbf{«Zнг»}. 

\begin{figure}[H]
\centering
\includegraphics[width=0.5\textwidth,height=0.5\textheight,keepaspectratio]{\fbpath/661.0031 ДЗ/img/_r.pdf}
\captionof{figure}{Характеристика срабатывания измерительного органа отстройки от нагрузки}\label{pict_dz:r}
\end{figure}

\item
Суммарная характеристика направленного РС, используемого в логике работы ступеней ДЗ приведена на рисунке \ref{pict_dz:rs}.  

\begin{figure}[H]
\centering
\includegraphics[width=0.75\textwidth,height=0.75\textheight,keepaspectratio]{\fbpath/661.0031 ДЗ/img/_rs.pdf}
\captionof{figure}{Характеристики срабатывания направленного РС}\label{pict_dz:rs}
\end{figure}

\item
Дополнительно для первой и третей ступени ДЗ предусмотрены направленные РС с охватом начала
координат, уставки по R, Х и φ1 которых совпадают с уставками ненаправленного РС соответствующей ступени.
Характеристика срабатывания направленного РС с охватом начала координат имеет форму параллелограмма,
смещенного вдоль направления металлического КЗ на линии в третий и четвертый квадранты на величину до -0,1$\cdot X_{\textsl{ср}}$. 

\begin{figure}[H]
\centering
\includegraphics[width=0.75\textwidth,height=0.75\textheight,keepaspectratio]{\fbpath/661.0031 ДЗ/img/_nachalo.pdf}
\captionof{figure}{Характеристика срабатывания направленного РС с охватом начала координат}\label{pict_dz:nachalo}
\end{figure}

\item
Для  обнаружения  однофазных  КЗ  и  пуска  ступеней  ДЗ-ФЗ,  контролирующих  контуры  «фаза-земля»,
предусмотрен  орган  определения  вида  повреждения  (ООВП),  основанный  на  совместном  применении
измерительного органа тока нулевой последовательности (ИО 3I0 ООВП) с торможением от одного из фазных
токов (средний из  $I_A , I_B , I_C$) и измерительного органа напряжения нулевой последовательности (ИО 3U0 ООВП).
Торможение осуществляется от тока той фазы, значение в которой является средним между максимальным и
минимальным значениями токов в остальных двух фазах. В таком случае при междуфазном КЗ на землю
торможение максимальном, а при однофазном КЗ –- минимально. 

\item
Характеристика  срабатывания  ИО  3I0  ООВП  с  торможением  представлена  на  рисунке  \ref{pict_dz:torm}.  

\begin{figure}[H]
\centering
\includegraphics[width=0.65\textwidth,height=0.65\textheight,keepaspectratio]{\fbpath/661.0031 ДЗ/img/_torm.pdf}
\captionof{figure}{Характеристика срабатывания ИО 3I0 ООВП}\label{pict_dz:torm}
\end{figure}

\item
Порог срабатывания горизонтального участка задается уставкой \textbf{«3I0ср ООВП»}. Наклонный участок
характеристики проходит через точку, равной 1,25$ \cdot I_{\textsl{ном}}$ , под углом к горизонтальной оси, соответствующим
коэффициенту торможения $k_T$. Коэффициент торможения задается уставкой \textbf{«Кт ООВП»}. При превышении
тормозного тока уставки \textbf{«Iср блок. торм.»} ООВП блокируется.

\begin{figure}[H]
\centering
\includegraphics[width=0.65\textwidth,height=0.65\textheight,keepaspectratio]{\fbpath/661.0031 ДЗ/img/OOVP.pdf}
\captionof{figure}{Функциональная схема ООВП}\label{pict_dz:OOVP}
\end{figure}

\item
Ввод  в  работу  ступеней  ДЗ,  работающих  по  контуру  «фаза-земля»  может  осуществляться  с
дополнительным  контролем  напряжения  нулевой  последовательности от ТН 1СШ  с помощью  программного
ключа \textbf{«Контр.  3U0 ООВП»}.  Напряжение  срабатывания  ИО  3U0  ООВП  задается  уставкой  \textbf{«3U0ср ООВП»}. ИО
3U0 ООВП может работать по напряжению $3U_{\textsl{0 звезда}}$, рассчитанному на основе фазных напряжений, или по
напряжению $3U_{\textsl{0 треуг.}}$   --  выбирается  программным  переключателем  \textbf{«Выбор  3U0»}, имеющим  два
соответствующих положения: \textsl{«Расчетн.»} и \textsl{«Измерен.»}.

\item
Переменная $3U_{\textsl{0 треуг.}}$ зависит от положения программного переключателя \textbf{«Схема подключения цепей
разомкнутого треугольника»}.  Данный  переключатель  позволяет  учесть  различные  схемы  соединения
дополнительной  обмотки  ТН.  Если  переключатель  выставлен  в  положение  «$U_{\textsl{НК}}$»,  то  переменной  $3U_{\textsl{0 треуг.}}$
присваивается  значение  напряжения,  измеренного  по  аналоговому  входу  «$U_{\textsl{ИФ}}(U_{\textsl{НК}})$».  Если  программный
переключатель {\textbf{«Схема подключения цепей треугольника»}} выставлен в положение «$U_{\textsl{НИ}}/U_{\textsl{ИК}}$», «$U_{\textsl{НФ}}/U_{\textsl{ФК}}$»
или  «$U_{\textsl{НИ}}/U_{\textsl{ИФ}}/U_{\textsl{ФК}}$», то переменная $3U_{\textsl{0 треуг.}}$ рассчитывается в соответствии с формулами, приведенными в
приложении \ref{app:signals}. Программный переключатель \textbf{«Выбор 3U0»} аналогично влияет на работу ОНМНП. 

\item
В составе логики дистанционной защиты предусмотрено четыре ступени от междуфазных замыканий с
контролем трех контуров «фаза-фаза». 

\vspace{2mm}
\begin{figure}[H]
\centering
\includegraphics[width=0.65\textwidth,height=0.65\textheight,keepaspectratio]{\fbpath/661.0031 ДЗ/img/DZ_FF_1.pdf}
\captionof{figure}{Функциональная схема ДЗ-ФФ-1(3)}\label{pict_dz:DZ_FF_1}
\end{figure}

\item
Первая и третья ступень ДЗ от междуфазных КЗ содержат две «подступени»: быстродействующая 
ДЗ-ФФ-1(б) и медленнодействующая ДЗ-ФФ-1(м), которая вводится программным ключом \textbf{«ДЗ-ФФ-1(3)(м)»}. 
Всю ступень можно ввести/вывести программным ключом \textbf{«ДЗ-ФФ-1(3)»}. Контроль от блокировки при качаниях
(БК) для каждой подступени задается отдельными программными переключателями \textbf{«БК ДЗ-ФФ-1(3)(б)»} и 
\textbf{«БК ДЗ-ФФ-1(3)(м)»}. Действие подступеней осуществляется с регулируемыми независимыми выдержками времени
\textbf{«Тср ДЗ-ФФ-1(3)(б)»} и \textbf{«Тср ДЗ-ФФ-1(3)(м)»}. Пороги срабатывания задаются уставками \textbf{«Хср-ФФ-1(3)»}, \textbf{«Rср-ФФ-1(3)»} и \textbf{«Ф1-ФФ-1(3)}.

\item
Предусмотрена возможность работы ДЗ-ФФ-1(3) с охватом начала координат, задается программным
ключем \textbf{«Режим охвата НК ДЗ-ФФ-1(3)»}. А также, предусмотрена возможность ввода подхвата
направленного пуска от ненаправленного пуска этих же ступеней, задается программным ключем \textbf{«Подхват
РС-ФФ-1(3)»}. Тем самым гарантируется правильное действие защиты при близких замыканиях. 

\vspace{2mm}
\begin{figure}[H]
\centering
\includegraphics[width=0.65\textwidth,height=0.65\textheight,keepaspectratio]{\fbpath/661.0031 ДЗ/img/DZ_FF_2.pdf}
\captionof{figure}{Функциональная схема ДЗ-ФФ-2(4)}\label{pict_dz:DZ_FF_1}
\end{figure}

\item
Вторая и четвертая ступень ДЗ от междуфазных КЗ вводится программным ключом \textbf{«ДЗ-ФФ-2(4)»}. 
Контроль от блокировки при качаниях (БК) задается программными переключателями \textbf{«БК ДЗ-ФФ-2(4)»}.
Действие ступеней осуществляется с регулируемыми независимыми выдержками времени \textbf{«Тср ДЗ-ФФ-2(4)»}.
Пороги срабатывания задаются уставками \textbf{«Хср-ФФ-2(4)»}, \textbf{«Rср-ФФ-2(4)»} и \textbf{«Ф1-ФФ-2(4)»}.

\item
Оперативный ввод/вывод ДЗ-ФФ производится с помощью БУ (блоков управления функциями) \textit{«Вывод ДЗ»}.

\item
В составе логики дистанционной защиты предусмотрено четыре ступеней от однофазных замыканий с
контролем трех контуров «фаза-земля».

\vspace{2mm}
\begin{figure}[H]
\centering
\includegraphics[width=0.65\textwidth,height=0.65\textheight,keepaspectratio]{\fbpath/661.0031 ДЗ/img/DZ_FZ_1.pdf}
\captionof{figure}{Функциональная схема ДЗ-ФЗ-1(3)}\label{pict_dz:DZ_FZ_1}
\end{figure}

\item
Первая и третья ступень ДЗ от однофазных КЗ содержат две «подступени»: быстродействующая 
ДЗ-ФЗ-1(б) и медленнодействующая ДЗ-ФЗ-1(м), которая вводится программным ключом \textbf{«ДЗ-ФЗ-1(3)(м)»}. 
Всю ступень можно ввести/вывести программным ключом \textbf{«ДЗ-ФЗ-1(3)»}. Контроль от блокировки при качаниях
(БК) для каждой подступени задается отдельными программными переключателями \textbf{«БК ДЗ-ФЗ-1(3)(б)»} и 
\textbf{«БК ДЗ-ФЗ-1(3)(м)»}. Действие подступеней осуществляется с регулируемыми независимыми выдержками
времени \textbf{«Тср ДЗ-ФЗ-1(3)(б)»} и \textbf{«Тср ДЗ-ФЗ-1(3)(м)»}. Пороги срабатывания задаются уставками \textbf{«Хср-ФЗ-
1(3)»}, \textbf{«Rср-ФЗ-1(3)»} и \textbf{«Ф1-ФЗ-1(3)»}.

\item
Предусмотрена возможность работы ДЗ-ФЗ-1(3) с охватом начала координат, задается программным
ключом \textbf{«Режим охвата НК ДЗ-ФЗ-1(3)»}. А также, предусмотрена возможность ввода подхвата направленного
пуска от ненаправленного пуска этих же ступеней, задается программным ключом \textbf{«Подхват РС-ФЗ-1(3)»}. Тем
самым гарантируется правильное действие защиты при близких замыканиях. 

\vspace{2mm}
\begin{figure}[H]
\centering
\includegraphics[width=0.65\textwidth,height=0.65\textheight,keepaspectratio]{\fbpath/661.0031 ДЗ/img/DZ_FZ_2.pdf}
\captionof{figure}{Функциональная схема ДЗ-ФЗ-2(4)}\label{pict_dz:DZ_FZ_2}
\end{figure}

\item
Вторая и четвертая ступень ДЗ от междуфазных КЗ вводится программным ключом \textbf{«ДЗ-ФЗ-2(4)»}. 
Контроль от блокировки при качаниях (БК) задается программными переключателями \textbf{«БК ДЗ-ФЗ-2(4)»}.
Действие ступеней осуществляется с регулируемыми независимыми выдержками времени \textbf{«Тср ДЗ-ФЗ-2(4)»}.
Пороги срабатывания задаются уставками \textbf{«Хср-ФЗ-2(4)»}, \textbf{«Rср-ФЗ-2(4)»} и \textbf{«Ф1-ФЗ-2(4)»}.

\item
Оперативный ввод/вывод ДЗ-ФЗ производится с помощью БУ (блоков управления функциями) \textit{«Вывод ДЗ»}. 

\item
Для исключения ложного действия при возникновении неисправности в цепях напряжения ступени ДЗ
могут быть заблокированы от БНН. Блокировка от БНН вводится отдельно для каждой ступени ДЗ с помощью
программных переключателей \textbf{«БНН ДЗ-ФФ-1(2…4)»} и \textbf{«БНН ДЗ-ФЗ-1(2…4)»}.

\item
Дистанционная защита контролируется блокировкой при качаниях, предотвращающей ложное
срабатывание ДЗ при качаниях и асинхронном ходе. В устройстве реализовано два способа блокировки:
\begin{itemize}
\item по приращению токов прямой и обратной последовательности;
\item по скорости изменения величины сопротивления dZ/dt.
\end{itemize}

\item
Контроль ступеней ДЗ от блокировки при качаниях вводится программными переключателями
\textbf{«БК ДЗ-ФФ(ФЗ)-1(2…4)»}, имеющими четыре положения:
\begin{itemize}
\item \textit{«Выведено»} -- ступень работает без контроля от БК;
\item \textit{«БК-I-б»} -- ступень пускается от сигнала ввода быстродействующих ступеней;
\item \textit{«БК-I-м»} -- ступень пускается от сигнала ввода медленнодействующих ступеней;
\item \textit{«БК-Z»} -- ступень пускается от сигнала БК по скорости изменения величины сопротивления.
\end{itemize}

\item
Если за время пуска от БК происходит срабатывание РС ступени ДЗ, то разращение от БК
подхватывается и удерживается до пропадания условий срабатывания соответствующей ступени.

\item
В составе логики дистанционной защиты предусмотрено два режима оперативного ускорения ДЗ.

\vspace{2mm}
\begin{figure}[H]
\centering
\includegraphics[width=1\textwidth,height=1\textheight,keepaspectratio]{\fbpath/661.0031 ДЗ/img/OA_DZ_without_T.pdf}
\captionof{figure}{Функциональная схема ОУ ДЗ без выдержки времени}\label{pict_dz:OA_DZ_without_T}
\end{figure}

\item
Оперативное ускорение ДЗ без выдержки времени предназначено для оперативного ускорения ступени
ДЗ-ФФ при выполнении переключений в сети. Ввод в работу ОУ-бв ДЗ и выбор ускоряемой ступени задается
программным переключателем \textbf{«ОУ-бв ДЗ»}. Предусмотрена возможность вывода направленности ускоряемой
ступени с помощью программного ключа \textbf{«Направл ДЗ при ОУ-бв»}. Действие ускорения осуществляется с
регулируемой независимой выдержкой времени \textbf{«Тср ОУ-бв ДЗ»}. Контроль ускорения от блокировки при
качаниях (БК) задается программным переключателем \textbf{«БК ОУ-бв ДЗ»}, а блокировка от БНН вводится с
помощью программного ключа \textbf{«БНН ОУ-бв ДЗ»}. 

\item
Оперативный ввод ОУ-бв ДЗ производится с помощью команды управления \textit{«Ввод ОУ-бв ДЗ»}.

\vspace{2mm}
\begin{figure}[H]
\centering
\includegraphics[width=1\textwidth,height=1\textheight,keepaspectratio]{\fbpath/661.0031 ДЗ/img/OA_DZ.pdf}
\captionof{figure}{Функциональная схема ОУ ДЗ}\label{pict_dz:OA_DZ}
\end{figure}

\item
Оперативное ускорение ДЗ предназначено для оперативного ускорения ступени ДЗ-ФФ при выводе ДЗШ.
Ввод в работу ОУ ДЗ и выбор ускоряемой ступени задается программным переключателем \textbf{«ОУ ДЗ»}.
Предусмотрена возможность вывода направленности ускоряемой ступени с помощью программного ключа
\textbf{«Направл ДЗ при ОУ»}. Действие ускорения осуществляется с регулируемой независимой выдержкой времени
\textbf{«Тср ОУ ДЗ»}. Контроль ускорения от блокировки при качаниях (БК) задается программным переключателем
\textbf{«БК ОУ ДЗ»}, а блокировка от БНН вводится с помощью программного ключа \textbf{«БНН ОУ ДЗ»}. 

\item
Оперативный ввод ОУ ДЗ производится с помощью команды управления \textit{«Ввод ОУ ДЗ»}.

\item
В составе логики дистанционной защиты предусмотрен пуск автоматического ускорения ступеней ДЗ.

\vspace{2mm}
\begin{figure}[H]
\centering
\includegraphics[width=1\textwidth,height=1\textheight,keepaspectratio]{\fbpath/661.0031 ДЗ/img/AA_DZ.pdf}
\captionof{figure}{Функциональная схема пуска АУ ДЗ}\label{pict_dz:AA_DZ}
\end{figure}

\item
Ввод в работу пуска АУ ДЗ и выбор ускоряемой ступени задается программным переключателем \textbf{«АУ
ДЗ»}. Предусмотрена возможность вывода направленности ускоряемой ступени с помощью программного
ключа \textbf{«Направл ДЗ при АУ»}. Контроль ускорения от блокировки при качаниях (БК) задается программным
переключателем \textbf{«БК АУ ДЗ»}, а блокировка от БНН вводится с помощью программного ключа \textbf{«БНН АУ ДЗ»}.

\small
\begin{longtable}{|>{\centering\arraybackslash}m{5.3cm}|>{\centering\arraybackslash}m{3.3cm}|>{\centering\arraybackslash}m{4.2cm}|>{\centering\arraybackslash}m{1.8cm}|>{\centering\arraybackslash}m{1cm}|}
% Первая страница: полный заголовок
\caption{Перечень уставочных параметров ДЗ\hfill\vspace{-0.5\baselineskip}}\label{dz:tbl1}\\
\hline
\rowcolor{gray!30}
Параметр & Условное обозначение на схеме & Значение/ Диапазон & Единица измерения & Шаг \\
\hline
\endfirsthead

% Промежуточные страницы: «Продолжение таблицы»
\caption*{\hspace{3pt}\emph{Продолжение таблицы \ref{dz:tbl1}}\hfill\ }\vspace{-0.5\baselineskip} \\
\hline
\rowcolor{gray!30}
Параметр & Условное обозначение на схеме & Значение/ Диапазон & Единица измерения & Шаг \\
\hline
\endhead

%===m>ДЗ|-
%===m
\end{longtable}
\normalsize

\end{enumerate}
